% Options for packages loaded elsewhere
\PassOptionsToPackage{unicode}{hyperref}
\PassOptionsToPackage{hyphens}{url}
\PassOptionsToPackage{dvipsnames,svgnames,x11names}{xcolor}
\PassOptionsToPackage{space}{xeCJK}
\documentclass[
  chinese,
  12pt,
  a4paper,
  oneside]{article}
\usepackage{xcolor}
\usepackage[top=24mm,bottom=26mm,left=24mm,right=24mm]{geometry}
\usepackage{amsmath,amssymb}
\setcounter{secnumdepth}{-\maxdimen} % remove section numbering
\usepackage{iftex}
\ifPDFTeX
  \usepackage[T1]{fontenc}
  \usepackage[utf8]{inputenc}
  \usepackage{textcomp} % provide euro and other symbols
\else % if luatex or xetex
  \usepackage{unicode-math} % this also loads fontspec
  \defaultfontfeatures{Scale=MatchLowercase}
  \defaultfontfeatures[\rmfamily]{Ligatures=TeX,Scale=1}
\fi
\usepackage{lmodern}
\ifPDFTeX\else
  % xetex/luatex font selection
  \setmonofont[]{Menlo}
  \ifXeTeX
    \usepackage{xeCJK}
    \setCJKmainfont[]{Songti TC}
  \fi
  \ifLuaTeX
    \usepackage[]{luatexja-fontspec}
    \setmainjfont[]{Songti TC}
  \fi
\fi
% Use upquote if available, for straight quotes in verbatim environments
\IfFileExists{upquote.sty}{\usepackage{upquote}}{}
\IfFileExists{microtype.sty}{% use microtype if available
  \usepackage[]{microtype}
  \UseMicrotypeSet[protrusion]{basicmath} % disable protrusion for tt fonts
}{}
\usepackage{setspace}
\makeatletter
\@ifundefined{KOMAClassName}{% if non-KOMA class
  \IfFileExists{parskip.sty}{%
    \usepackage{parskip}
  }{% else
    \setlength{\parindent}{0pt}
    \setlength{\parskip}{6pt plus 2pt minus 1pt}}
}{% if KOMA class
  \KOMAoptions{parskip=half}}
\makeatother
\ifLuaTeX
\usepackage[bidi=basic,shorthands=off,provide=*]{babel}
\else
\usepackage[bidi=default,shorthands=off,provide=*]{babel}
\fi
\ifLuaTeX
  \usepackage{selnolig} % disable illegal ligatures
\fi
\setlength{\emergencystretch}{3em} % prevent overfull lines
\providecommand{\tightlist}{%
  \setlength{\itemsep}{0pt}\setlength{\parskip}{0pt}}
\usepackage{xcolor}
\usepackage{fancyhdr}
\usepackage{setspace}

\setstretch{1.18}
\setlength{\parindent}{2em}
\setlength{\parskip}{0.45em}
\setlength{\headheight}{14pt}

\pagestyle{fancy}
\fancyhf{}
\fancyhead[L]{惠文國二下數學 PBL 自學教材}
\fancyhead[R]{\thepage}
\renewcommand{\headrulewidth}{0.2pt}

\AtBeginDocument{
  \hypersetup{
    unicode=true,
    pdfborder={0 0 0}
  }
}
\usepackage{bookmark}
\IfFileExists{xurl.sty}{\usepackage{xurl}}{} % add URL line breaks if available
\urlstyle{same}
\hypersetup{
  pdftitle={惠文國二下數學段考對齊讀書計畫},
  pdflang={zh-TW},
  colorlinks=true,
  linkcolor={blue},
  filecolor={Maroon},
  citecolor={Blue},
  urlcolor={blue},
  pdfcreator={LaTeX via pandoc}}

\title{惠文國二下數學段考對齊讀書計畫}
\author{}
\date{}

\begin{document}
\maketitle

{
\setcounter{tocdepth}{2}
\tableofcontents
}
\setstretch{1.18}
版本日期：2026-03-08\\
整理依據：惠文高中公開課程計畫 \texttt{表三彙整0813(8年級).pdf}\\
用途：給自學生快速對照段考範圍與建議讀法

\section{一、先看結論}\label{ux4e00ux5148ux770bux7d50ux8ad6}

\subsection{第一次定期評量前}\label{ux7b2cux4e00ux6b21ux5b9aux671fux8a55ux91cfux524d}

\begin{itemize}
\tightlist
\item
  建議完成：

  \begin{itemize}
  \tightlist
  \item
    \texttt{1-1} 認識數列與等差數列
  \item
    \texttt{1-2} 等差級數
  \item
    \texttt{1-3} 等比數列
  \end{itemize}
\item
  對應教材：

  \begin{itemize}
  \tightlist
  \item
    \texttt{單元\ 1-1、1-2　四大島嶼的節奏密碼}
  \item
    \texttt{單元\ 1-3　倍率傳訊與等比數列}
  \end{itemize}
\end{itemize}

\subsection{第二次定期評量前}\label{ux7b2cux4e8cux6b21ux5b9aux671fux8a55ux91cfux524d}

\begin{itemize}
\tightlist
\item
  建議完成：

  \begin{itemize}
  \tightlist
  \item
    \texttt{2-1} 變數與函數
  \item
    \texttt{2-2} 線型函數與圖形
  \item
    \texttt{3-1} 內角與外角
  \item
    \texttt{3-2} 尺規作圖與三角形的全等
  \end{itemize}
\item
  對應教材：

  \begin{itemize}
  \tightlist
  \item
    \texttt{單元\ 2-1、2-2　水安宮站的成本函數}
  \item
    \texttt{單元\ 3-1　角度偵查與內角外角}
  \item
    \texttt{單元\ 3-2　尺規作圖與公平服務點}
  \end{itemize}
\end{itemize}

\subsection{第三次定期評量前}\label{ux7b2cux4e09ux6b21ux5b9aux671fux8a55ux91cfux524d}

\begin{itemize}
\tightlist
\item
  建議完成：

  \begin{itemize}
  \tightlist
  \item
    \texttt{3-3} 全等三角形的應用
  \item
    \texttt{3-4} 中垂線與角平分線性質
  \item
    \texttt{3-5} 三角形的邊角關係
  \item
    \texttt{4-1} 平行線與截角性質
  \item
    \texttt{4-2} 平行四邊形
  \item
    \texttt{4-3} 特殊四邊形與梯形
  \end{itemize}
\item
  對應教材：

  \begin{itemize}
  \tightlist
  \item
    \texttt{單元\ 3-3　全等三角形與穩定支架}
  \item
    \texttt{單元\ 3-4　中垂線與角平分線指令}
  \item
    \texttt{單元\ 3-5　三角形的邊角關係}
  \item
    \texttt{單元\ 4-1　平行線與截角暗號}
  \item
    \texttt{單元\ 4-2　平行四邊形與展區版面}
  \item
    \texttt{單元\ 4-3　特殊四邊形與梯形總檢}
  \end{itemize}
\end{itemize}

\section{二、每週怎麼讀}\label{ux4e8cux6bcfux9031ux600eux9ebcux8b80}

\subsection{讀新單元的順序}\label{ux8b80ux65b0ux55aeux5143ux7684ux9806ux5e8f}

\begin{enumerate}
\def\labelenumi{\arabic{enumi}.}
\tightlist
\item
  先聽 \texttt{voice}
\item
  再看直覺問題
\item
  再看公式與手算
\item
  做 \texttt{checkpoint}
\item
  做章末練習
\item
  最後口頭解釋一次
\end{enumerate}

\subsection{一次讀書的最小單位}\label{ux4e00ux6b21ux8b80ux66f8ux7684ux6700ux5c0fux55aeux4f4d}

\begin{itemize}
\tightlist
\item
  \texttt{25} 分鐘看新內容
\item
  \texttt{10} 分鐘手算
\item
  \texttt{5} 分鐘口頭重講
\end{itemize}

\subsection{段考前最後一週}\label{ux6bb5ux8003ux524dux6700ux5f8cux4e00ux9031}

\begin{itemize}
\tightlist
\item
  第 \texttt{7} 天到第 \texttt{5} 天：重做所有例題
\item
  第 \texttt{4} 天到第 \texttt{3} 天：重做章末練習
\item
  第 \texttt{2} 天：整理最容易錯的公式和判斷點
\item
  第 \texttt{1} 天：只做輕量複習，不再硬塞新題
\end{itemize}

\section{三、這份計畫怎麼用}\label{ux4e09ux9019ux4efdux8a08ux756bux600eux9ebcux7528}

\begin{itemize}
\tightlist
\item
  如果你跟著學校正常上課，就拿它來對照自己有沒有落後。
\item
  如果你是自主超前，就拿它來安排每次段考前的回頭複習。
\item
  如果老師公布的實際考試範圍和公開課程計畫不同，以老師現場宣布為準。
\end{itemize}

\end{document}
